\documentclass[11pt]{article}
\usepackage[a4paper,margin=27mm]{geometry}
\usepackage{amsmath,amssymb,amsthm,booktabs,hyperref}
\usepackage{xurl}
\hypersetup{colorlinks=true,linkcolor=blue,urlcolor=blue,citecolor=blue}
\newtheorem{theorem}{Theorem}
\newtheorem{lemma}{Lemma}
\newtheorem{corollary}{Corollary}
\newcommand{\RR}{\mathbb R}
\newcommand{\dist}{\operatorname{dist}}
\title{Globally optimal quadratic deficit bounds for two- and three-point reciprocal products}
\author{Mohith (MokiMeow)}
\date{11 October 2026\\\small Formal proof note; not peer reviewed}
\begin{document}
\maketitle
\begin{abstract}
We give machine-checked proofs of quadratic error bounds with optimal global coefficients for
the product of a sum and a reciprocal sum on a positive interval. In the
two-point weighted problem, the optimal coefficient for squared Euclidean
distance to the two maximizing endpoint pairs is $pq/(mM)$. In the
unweighted three-point problem, the optimal coefficient for distance to
the six endpoint triples with one or two upper endpoints is $2/(mM)$. The three-point proof uses
six positive-polynomial certificates and an explicit covering of the
ordered domain. A Lean development also verifies the classical exact
fixed-weight partition formula as background. All formal theorems are
proved without unfinished proof placeholders or additional axioms.
The present note records verified statements and proofs; first-publication
priority and editorial novelty are not claimed.
\end{abstract}

\section{Setting and relation to classical results}
Throughout, $0<m<M$ and every input belongs to $[m,M]$. The classical
Kantorovich inequality bounds
\[
 \left(\sum_i w_i x_i\right)\left(\sum_i w_i/x_i\right)
 \le \frac{(m+M)^2}{4mM}\left(\sum_i w_i\right)^2,
 \qquad w_i\ge0.
\]
Newman gives an elementary scalar reduction and convexity proof
\cite{newman}. The exact fixed-weight maximum is also established
literature: if $W=\sum_i w_i$ and
$\delta^2=\min_S(\sum_{i\in S}w_i-\sum_{i\notin S}w_i)^2$, then
\[
 Q_* = \frac{(m+M)^2 W^2-(M-m)^2\delta^2}{4mM}.
\]
This formula appears in Niculescu--Persson, Section 1.4, Exercise 1
\cite{niculescu}. For equal weights, Lupa\c{s} gives the even/odd correction
\cite{lupas}. We use these results as historical context and supply
independent formal proofs; they are not presented as new mathematics.

The question here is quantitative: if an input's product is close to its
maximum, how close is that input to a maximizing endpoint configuration?
The distance used below is the ordinary unweighted Euclidean distance on
the coordinate space. Changing that metric changes the coefficient.

\section{Two weighted points}
For $p,q>0$, put
\[
 Q(x,y)=(px+qy)(p/x+q/y),\qquad
 Q_*=(p+q)^2+\frac{pq(M-m)^2}{mM}.
\]
The maximizing set is $E_2=\{(m,M),(M,m)\}$. Define
\[
 d_2^2(x,y)=\min\{(x-m)^2+(y-M)^2,
                         (x-M)^2+(y-m)^2\}.
\]
\begin{theorem}\label{thm:two}
For every $x,y\in[m,M]$,
\[
 Q_*-Q(x,y)\ \ge\ \frac{pq}{mM}\,d_2^2(x,y).
\]
The coefficient $pq/(mM)$ is the largest coefficient valid for every
input in the interval. Moreover, $Q(x,y)=Q_*$ if and only if
$(x,y)\in E_2$.
\end{theorem}
\begin{proof}
Expansion gives the exact deficit identity
\[
 Q_*-Q=\frac{pq}{mM}
 \left((M-m)^2-\frac{mM(x-y)^2}{xy}\right).
\]
By symmetry assume $x\le y$, so
$d_2^2=(x-m)^2+(M-y)^2$. The numerator of the required scalar inequality
has the following decomposition:
\begin{align*}
 &xy\big((M-m)^2-(x-m)^2-(M-y)^2\big)-mM(y-x)^2\\
 &\quad=(x-m)(y-x)^2y+2xy(x-m)(y-x)\\
 &\qquad\quad+(M-y)(y-x)\big(2x^2+(2x-m)(y-x)\big)\\
 &\qquad\quad+2xy(x-m)(M-y).
\end{align*}
Every term is nonnegative. Divide by $xy>0$ and use the deficit identity.
At $x=y=m$, the distance equals $(M-m)^2>0$ and the deficit equals
$pq(M-m)^2/(mM)$, proving optimality. The bound forces $d_2^2=0$ for a
maximizer, and direct evaluation verifies both endpoint pairs.
\end{proof}

\begin{corollary}
If $Q_*-Q(x,y)\le\varepsilon$, then
$d_2^2(x,y)\le mM\varepsilon/(pq)$.
\end{corollary}
The inequality itself permits $p=0$ or $q=0$ and $m=M$. The extremizer
classification and the distance corollary above require positive weights.
The formal optimal-coefficient statement includes zero weights, in which
case the optimal coefficient is zero.

Optimality concerns the single coefficient valid on the entire interval.
We do not assert that squared distance gives the optimal local rate near an
endpoint maximizer; the distance corollary can be loose in that regime.

\section{Three unweighted points}
Put
\[
 Q(x,y,z)=(x+y+z)(1/x+1/y+1/z),\qquad
 Q_*=9+\frac{2(M-m)^2}{mM}.
\]
Let $E_3$ consist of the six triples with coordinates in $\{m,M\}$ and
exactly one or exactly two coordinates equal to $M$. Write
$d_3^2=\dist((x,y,z),E_3)^2$, explicitly the minimum of the six sums of
squared coordinate differences.

\begin{theorem}\label{thm:three}
For every $x,y,z\in[m,M]$,
\[
 Q_*-Q(x,y,z)\ \ge\ \frac{2}{mM}\,d_3^2(x,y,z).
\]
The coefficient $2/(mM)$ is the largest coefficient valid for every
input in the interval.
\end{theorem}

\subsection{Complete polynomial-certificate proof}
Sort the inputs so that $m\le x\le y\le z\le M$. Define
\begin{align*}
 N_L={}&xyz\big(2(M-m)^2-2((x-m)^2+(y-m)^2+(M-z)^2)\big)\\
       &-mM\big((x-y)^2z+(x-z)^2y+(y-z)^2x\big).
\end{align*}
Define $N_U$ by replacing $(y-m)^2$ by $(y-M)^2$.
Exact algebra gives
\[
 Q_*-Q-\frac{2}{mM}\big((x-m)^2+(y-m)^2+(M-z)^2\big)
   =\frac{N_L}{mMxyz},
\]
and the analogous identity for $N_U$. Both denominators are positive.
It is therefore sufficient to prove $N_L\ge0$ when $2y\le m+M$ and
$N_U\ge0$ when $2y\ge m+M$.

Here is an exhaustive cover of the first region. In each row, all listed
parameters are nonnegative; the columns give $x,y,z,M$ respectively.
\begin{center}\small
\begin{tabular}{@{}cllll@{}}
\toprule
 & $x$ & $y$ & $z$ & $M$\\
\midrule
0 & $m+a$ & $m+a+t$ & $m+2a+2t+r$ & $m+2a+2t+r+b$\\
1 & $m+a$ & $m+a+v+t$ & $m+2a+2v+t$ & $m+2a+2v+2t+r$\\
2 & $m+u+a$ & $m+u+a+t$ & $m+2u+a+t$ & $m+2u+2a+2t+r$\\
\bottomrule
\end{tabular}
\end{center}
To verify coverage, first test $2y\le z+m$. If true, use row 0 with
$(a,t,r,b)=(x-m,y-x,z-2y+m,M-z)$. Otherwise test $x-m\le z-y$.
If true, use row 1 with
\[
 (a,v,t,r)=(x-m,z-y-(x-m),2y-z-m,M+m-2y).
\]
If false, use row 2 with
\[
 (u,a,t,r)=(z-y,x-m-(z-y),y-x,M+m-2y).
\]
The interval, ordering, and half-region inequalities prove every parameter
nonnegative. Substitution reproduces the original four coordinates exactly.

Substituting these three rows into $N_L$ gives positive-coefficient
polynomials with respectively 66, 71, and 83 monomials. Their complete
integer-coefficient expressions are supplied in the formal supplement.
For the upper half, reflect the ordered triple by
\[
 (x,y,z)\longmapsto(m+M-z,m+M-y,m+M-x).
\]
Substituting the same three rows into the reflected $N_U$ gives
positive-coefficient polynomials with 89, 107, and 115 monomials.
Reflection maps the upper half to the lower half, so this also exhausts
the upper region. Each of the six expansion identities and its positivity
is checked separately in Lean.

The squared distance $d_3^2$ is no larger than the corresponding squared
distance used in $N_L$ or $N_U$. This proves the theorem for sorted inputs.
The product and the six-point distance are permutation invariant; the
formal proof checks all six orderings, completing the unrestricted case.

Finally, at $x=y=z=m$, the distance is $(M-m)^2$ and the deficit is
$2(M-m)^2/(mM)$. Thus a larger coefficient is impossible. Endpoint
evaluation gives $Q(m,m,M)=Q_*$, confirming the asserted maximum.

\section{Verification and scope}
The accompanying source pins Lean 4.34.0 and mathlib commit
\begin{center}\small\texttt{5ed2965256430c3649e86755f9576b54eca72435}\end{center}
The formal source proves Theorems \ref{thm:two} and \ref{thm:three},
characterizes their best coefficients, and gives the two-point distance
certificate. The polynomial generator helps discover the certificates.
The formal checker verifies the resulting identities and their signs;
it does not trust an external numerical oracle.

Every declared result in the verification record has a complete proof.
Axiom audits permit only the standard Lean logical axioms
\texttt{propext}, \texttt{Classical.choice}, and \texttt{Quot.sound}.
Targeted kernel replay checks the saved declarations. This is not a claim
of verification in a second independent proof assistant, physical-system
validation, or floating-point numerical error certification.

Exact finite computations suggest a dimension-dependent extension for
larger unweighted systems. That extension is not proved here and is not
part of the publication's verified claims.

\subsection{Exact examples}
For $m=1,M=4$ and unit weights, the following calculations are also
checked in the formal source. A stability equality can have positive
distance, so it must not be confused with an objective maximizer.
\begin{center}
\begin{tabular}{@{}llrr@{}}
\toprule
Problem & Input & Squared distance & Deficit\\
\midrule
Two points & $(1,1)$ & $9$ & $9/4$\\
Three points & $(1,1,1)$ & $9$ & $9/2$\\
Three points & $(1,5/2,4)$ & $9/4$ & $9/8$\\
\bottomrule
\end{tabular}
\end{center}

\section{Prior-work and publication status}
The fixed-weight partition formula and the equal-weight parity correction
are explicitly attributed above. A recent weighted Schweitzer refinement
\cite{alqahtani} is relevant neighboring work; its accessible HTML omits
the equations and its full remainder could not be compared. Consequently,
absence of an exact search match is not evidence of priority for our
distance estimates. The contribution of this release is a reproducible,
fully checked development with explicit globally optimal coefficients and finite
certificates. Editorial acceptance and first-publication novelty have not
been established.

\begingroup\small\raggedright
\begin{thebibliography}{9}
\bibitem{newman} M. Newman, \emph{Kantorovich's inequality},
Journal of Research of the National Bureau of Standards B 64B (1960),
33--34. \href{https://nvlpubs.nist.gov/nistpubs/jres/64B/jresv64Bn1p33_A1b.pdf}{Original paper (NIST PDF)}.
\bibitem{niculescu} C. P. Niculescu and L.-E. Persson,
\emph{Convex Functions and Their Applications}, manuscript dated
16 September 2004, Section 1.4, Exercise 1, p. 32.
\href{https://carmamaths.org/resources/jon/Preprints/Books/CUP/CUPold/np-convex.pdf}{Original manuscript (PDF)}.
\bibitem{lupas} A. Lupa\c{s}, \emph{A remark on the Schweitzer and
Kantorovich inequalities}, Publikacije Elektrotehni\v{c}kog Fakulteta,
Serija Matematika i Fizika (1972), 13--15.
\href{https://ictp.acad.ro/wp-content/uploads/2025/11/1972-Lupas-Publikacije-A-remark-on-the-Schweitzer-and-Kantorovich-inequalities.pdf}{Source PDF}.
\bibitem{alqahtani} Alqahtani and M. Z. Sarikaya,
\emph{Refined Weighted Schweitzer-Type Inequalities With Applications
to Statistical Moments}, Journal of Mathematics (2026), article 4822611.
\href{https://doi.org/10.1155/jom/4822611}{doi:10.1155/jom/4822611}.
\end{thebibliography}
\endgroup
\end{document}
